\BourbakiExercisesSection

\begin{enumerate}
\def\labelenumi{\arabic{enumi}.}
\tightlist
\item
  若 A 和 B 是 K-代数， \(\alpha\) 和 \(\beta\) 是从 A 到 B 中的两个 K-同态，则 A 到 B 中的一个 \((\alpha, \beta)\) -导子是指满足关系式
\end{enumerate}

\[
d (x y) = (d x) \alpha (y) + \beta (x) (d y) \quad \text { for } \quad x, y \quad \text { in } \quad A,
\]

换言之，即定义 1（第 2 号）意义下的导子，其中 \(\mathbf{A} = \mathbf{A}' = \mathbf{A}''\) , \(\mathbf{B} = \mathbf{B}' = \mathbf{B}''\) , \(d = d' = d''\) , \(\mu: \mathbf{A} \times \mathbf{A} \to \mathbf{A}\) 是乘法， \(\lambda_1: \mathbf{B} \times \mathbf{A} \to \mathbf{B}\) 是 K-双线性映射 \((z, x) \mapsto z\alpha(x)\) ，而 \(\lambda_2: \mathbf{A} \times \mathbf{B} \to \mathbf{B}\) 是 K-双线性映射 \((x, z) \mapsto \beta(x)z\) 。

从此以后，设 A 和 B 为交换域，且 \(\alpha \neq \beta\) 。证明此时存在一个元素 \(b \in \mathbf{B}\) ，使得 \(d\) 是 \((\alpha, \beta)\) -导子 \(x \mapsto b\alpha(x) - \beta(x)b\) 。（可以假设 \(d \neq 0\) 。首先证明 \(d\) 的核 N 是由 \(x \in \mathbf{A}\) （满足 \(\alpha(x) = \beta(x)\) ）组成的集合，然后证明对于 \(x \notin \mathbf{N}\) ，元素 \(b = d(x)/( \alpha(x) - \beta(x))\) 与 \(x\) 的选取无关。）

\begin{enumerate}
\def\labelenumi{\arabic{enumi}.}
\setcounter{enumi}{1}
\tightlist
\item
  设 K 是一个特征 \(\neq2\) 的交换域，A 和 B 为两个 K-代数，且 \(\alpha\) 为从 A 到 B 中的一个 K-线性映射；A 到 B 中的一个 \(\alpha\) -导子是指从 A 到 B 中的一个 K-线性映射 d，它满足关系式
\end{enumerate}

\[
d (x y) = (d x) \alpha (y) + \alpha (x) (d y) \quad \text { for } \quad x, y \quad \text { in } \quad A,
\]

换言之，即定义 1 意义下的一个导子，其中 \(\mathbf{A} = \mathbf{A}' = \mathbf{A}''\) , \(\mathbf{B} = \mathbf{B}' = \mathbf{B}''\) , \(d = d' = d''\) , \(\mu\) 是 \(\mathbf{A}\) 中的乘法， \(\lambda_1\) 是映射 \((z,x) \mapsto z\alpha(x)\) （从 \(\mathbf{B} \times \mathbf{A}\) 到 \(\mathbf{B}\) 中），而 \(\lambda_2\) 是映射 \((x,z) \mapsto \alpha(x)z\) （从 \(\mathbf{A} \times \mathbf{B}\) 到 \(\mathbf{B}\) 中）。

证明：若 B 是交换域、K 的一个扩域，且 \(d \neq 0\) ，则存在 B 中的 \(b \neq 0\) ，使得

\[
\alpha (x y) = \alpha (x) \alpha (y) + b (d x) (d y) \quad \text { for } \quad x, y \quad \text { in } \quad A.
\]

若 \(b = c^2\) ，其中 \(c \in \mathbf{B}\) , \(c \neq 0\) ，则存在两个同态 \(f, g\) （从 \(\mathbf{A}\) 到 \(\mathbf{B}\) 中），使得 \(\alpha(x) = (f(x) + g(x)) / 2\) , \(dx = (f(x) - g(x)) / 2c\) 。否则，存在一个二次代数（§ 2，第 3 号） \(\mathbf{B}'\) （在 \(\mathbf{B}\) 上），一个 \(\mu \in \mathbf{B}'\) ，使得 \(\mu^2 = c\) ，以及一个同态 \(f\) （从 \(\mathbf{A}\) 到 \(\mathbf{B}'\) 中），使得 \(\alpha(x) = (f(x) + \overline{f(x)}) / 2\) , \(dx = (f(x) - \overline{f(x)}) / 2\) 。

¶ 3. 设 K 为一个域，无论交换与否，E 为左向量空间 \(\mathrm{K}_{s}^{(\mathbf{N})}\) ，并设 \((e_{n})_{n\in\mathbf{N}}\) 为该空间的典范基。另一方面，设 \(\sigma\) 是 K 的一个自同态，且 d 是 K 到自身的一个 \((\sigma,1_{\mathbf{K}})\) -导子（习题 1），换言之，即 K 到自身的一个加法映射，使得

\[
d (\xi \eta) = (d \xi) \sigma (\eta) + \xi (d \eta).
\]

E 的元素是形如 \(\alpha_{n}e_{n}\) 的乘积的线性组合，并且在 E 上定义了一个乘法（E × E 到 E 中的 Z-双线性映射），其条件为

\[
\begin{array}{l} (\alpha e _ {n}) (\beta e _ {m}) = (\alpha e _ {n - 1}) (\sigma (\beta) e _ {m + 1} + (d \beta) e _ {m}) \text { for } n \geqslant 1 \\ (\alpha e _ {0}) (\beta e _ {m}) = (\alpha \beta) e _ {m}. \end{array}
\]

因此，在 E 上定义了一个（一般是非交换的）环结构，以 \(e_{0}\) 为单位元（与 K 的元素 l 等同）。若记 \(X = e_{1}\) ，则 \(e_{n} = X^{n}\) 对于 \(n \geqslant 1\) ，从而 E 的元素可以写成

\[
\alpha_ {0} + \alpha_ {1} X + \dots + \alpha_ {n} X ^ {n},
\]

根据乘法规则

\[
\mathrm{X} \alpha = \sigma (\alpha) \mathrm{X} + (d \alpha) \quad \text { for } \quad \alpha \in \mathrm{K}.
\]

这个环记作 K{[}X; σ, d{]} ，并称为关于 σ 和 d 的、系数在 K 中的 X 的非交换多项式环。~对于 E 中的一个非零元素 \(z = \sum_{j} \alpha_{j} X^{j}\) ，次数 deg(z) 是使得 \(\alpha_{j} \neq 0\) 的最大整数 j。

\begin{enumerate}
\def\labelenumi{(\alph{enumi})}
\tightlist
\item
  证明：如果 \(u, v\) 是两个元素 \(\neq 0\) （属于 \(E\) ），那么：
\end{enumerate}

\[
\begin{array}{r l} u v \neq 0 & \text { and } \quad \deg (u v) = \deg (u) + \deg (v); \\ \deg (u + v) & \leqslant \sup (\deg (u), \deg (v)) \quad \text { if } u + v \neq 0. \end{array}
\]

此外，若 \(\deg(u) \geqslant \deg(v)\) ，则存在 E 中的两个元素 w, r ，使得 \(u = wv + r\) ，并且要么 r = 0，要么 \(\deg(r) < \deg(v)\) （E 中的“欧几里得除法”）。

\begin{enumerate}
\def\labelenumi{(\alph{enumi})}
\setcounter{enumi}{1}
\tightlist
\item
  反之，设 R 是一个环，且存在一个映射 \(u \mapsto \deg(u)\) ，它把集合 \(R'\) （由 R 中 \(\neq 0\) 的元素组成）映入 N，并满足 (a) 的条件。证明由 0 与非零元素 \(u \in R\) （满足 \(\deg(u) = 0\) ）组成的集合 K 是一个（未必交换的）域。若 \(K \neq R\) 且 \(x\) 是 R 中的一个元素，使得整数 \(\deg(x) > 0\) 为尽可能小的，证明
\end{enumerate}

R 的每个元素都可以唯一地写成形式 \(\sum_{i} \alpha_j x^j\) ，其中 \(\alpha_j \in K\) 。（要看出 R 的每个元素都具有这种形式，可用反证法论证，考虑一个 \(u \in R\) （不具有这种形式的元素），且对它而言 \(\deg(u)\) 是尽可能小的。）证明存在 K 的一个自同态 \(\sigma\) 和 K 的一个 \((\sigma, 1_K)\) -导子 \(d\) ，使得 \(x\alpha = \sigma(\alpha)x + d\alpha\) 对所有 \(\alpha \in K\) 成立。由此推出 R 同构于一个域或一个环 K{[}X; \(\sigma, d]\) 。

\begin{enumerate}
\def\labelenumi{\arabic{enumi}.}
\setcounter{enumi}{3}
\tightlist
\item
  设 M 是类型为 N 的分次 K-模；证明每个次数为 r 的分次自同态都可以唯一地延拓为泛反交换代数 \(\mathbf{G}(\mathbf{M})\) （§ 7，习题 13）的次数为 r 的一个导子。
\end{enumerate}

* 5. 设 K 是特征为 \(p > 0\) 的域，B 是单不定元的有理函数域 K(X)，而 A 是 B 的子域 K(X \(^{p}\) ) 。证明典范同态 \(\Omega_{\mathbf{K}}(\mathbf{A}) \otimes_{\mathbf{A}} \mathbf{B} \to \Omega_{\mathbf{K}}(\mathbf{B})\) 不是单射。*

